\section{研究对象与决策问题}
\label{sec:01_01}

一项候选任务可能同时需要多种不可分资源，而其中任何一种资源也可能被其他任务需要。接受该任务，就意味着同时承诺它所要求的全部资源；若某项资源不能重复分配，这一选择便会排除其他需要该资源的候选。纯组合模型首先研究这种共同占用引起的选择限制。设$V$为有限资源集合，$E$为有限候选任务集合，每个候选任务$e\in E$以其所需资源的非空子集$e\subseteq V$表示，由此得到超图$H=(V,E)$：顶点表示资源，超边表示候选任务。对每个$e\in E$，无量纲二元变量$x_e\in\{0,1\}$表示是否接受任务，固定正权重$w_e\in\mathbb{R}_{>0}$表示统一评价尺度下的收益，本节将其视为无量纲分值；选择的目标是最大化总收益$\sum_{e\in E}w_ex_e$。资源$v\in V$的容量$b_v\in\mathbb{Z}_{\geq 0}$表示允许接受的占用份数，为无量纲计数。当前全部组合实验采用$b_v=1$，因此任意两个不同候选$e,f\in E$满足
\begin{equation}
 e\cap f\ne\varnothing,
 \quad x_e=1
 \quad\Longrightarrow\quad x_f=0.
 \label{s0101:shared_resource}
\end{equation}
这里的占用针对预先定义的候选及资源关系，不自动描述设备完成任务后如何释放、移动和再次使用。若同一设备可以在不同时段服务多个任务，就需要进一步描述占用时段和状态接续，不能把静态容量一解释为一台设备在整个班次内只能执行一次任务。

电力巡检将这种资源选择放入具体的空间与时间过程。巡检对象是杆塔及其需要观测的部位，执行资源包括车辆、无人机和具有独立身份的物理电池对。分别以有限集合$\mathcal{G}$、$\mathcal{V}$、$\mathcal{U}$和$\mathcal{B}$表示杆塔、车辆、无人机和物理电池对，其中车辆集合$\mathcal{V}$与组合模型的资源集合$V$区别使用；对每座杆塔$g\in\mathcal{G}$，以有限集合$\mathcal{I}_g$表示其必需观测项。车辆携带无人机与电池沿道路移动，无人机从满足条件的位置起飞，对杆塔不同部位实施观测，再与车辆会合回收；回收后的电池可以按其状态继续使用、换装或充电。因此，联合决策既要回答哪些观测交给哪架无人机，也要确定哪些观测组成同一航次、按什么顺序执行、由哪辆车支持、在哪里起飞和回收，以及使用哪一对电池。电池对具有实际位置、剩余能量和可用时刻，不能仅以一个可任意分配的总电量代替。

上述决定通过路径、事件时刻和资源状态相互影响。改变车辆停靠点，会改变无人机往返观测位置的距离，并影响车辆到达后续回收地点的时间；在同一航次中增加观测项，可能减少起降次数，也会改变飞行历时、耗电量与回收时刻；电池较晚返回车辆，则会推迟它开始充电或再次装机的时刻，进而限制后续航次。无人机到达杆塔附近也不等于完成观测：观测位置、姿态、可见性和成像质量须满足要求，连接这些观测状态的连续路径还须满足运动、障碍物避让、安全距离及能量条件。即使车辆路径和无人机路径分别可行，两者也可能因无法按时会合而不能共同执行。因此，任务分配必须与连续路径、观测状态、车辆移动、起降回收和电池周转一起决定。一个分配选择只有能够由相应的物理过程实现，才构成可执行的巡检安排。

巡检首先评价完整杆塔的完成数量，而不是累计拍摄次数。当前任务定义中，每座杆塔有九项必需观测，即对所有$g\in\mathcal{G}$均有$|\mathcal{I}_g|=9$；九项全部满足验收条件，才计为一座完整杆塔。缺少其中任何一项合格观测，该塔就不能计入完整完成数量，不同杆塔的零散观测也不能相互补足。用无量纲二元变量$z_g\in\{0,1\}$表示杆塔$g$在模型中的完整完成状态，其值为一当且仅当该塔全部必需观测均完成并满足模型规定的验收条件。相应的未完成杆塔数为
\begin{equation}
 N_{\mathrm{miss}}=\sum_{g\in\mathcal{G}}(1-z_g),
 \qquad
 N_{\mathrm{miss}}\in\{0,1,\ldots,|\mathcal{G}|\}.
 \label{s0101:uncompleted_towers}
\end{equation}
该量以杆塔座数计，整塔状态与各项观测之间的正式线性约束在第5章建立。以未完成数量为目标，允许在资源和物理条件限制下出现部分完成，并未预先要求所有杆塔必须完成。这里的完成状态是模型对计划的判定；执行后的实际验收仍须依据获得的观测资料和执行记录确认，不能由计划中的$z_g=1$直接替代。

在完整杆塔数量之后，还需比较全队结束任务的时间和无人机的机载耗电。令$T_{\mathrm{return}}\in\mathbb{R}_{\geq 0}$表示从统一任务起点到全队完成必要回收并返库的历时，单位为$\mathrm{s}$；令$E_{\mathrm{air}}\in\mathbb{R}_{\geq 0}$表示本次规划计入的机载耗电量，单位为$\mathrm{Wh}$。对于给定信息与给定情形下的一次规划，在满足物理及资源条件的方案中采用词典序目标
\begin{equation}
 \operatorname{lexmin}
 \bigl(N_{\mathrm{miss}},\ T_{\mathrm{return}},\ E_{\mathrm{air}}\bigr).
 \label{s0101:lexicographic_objective}
\end{equation}
比较两个方案时，先选择未完成杆塔较少的方案；仅当未完成数量相同，才比较全队返库历时；仅当前两项均相同，才比较机载耗电量。因此，少完成一座完整杆塔不能由更早返库或更低耗电补偿。一般加权混合目标通过权重规定不同指标之间的交换比例，可能接受完成数量与时间、耗电之间的交换，且权重须处理不同量纲；词典序则直接规定比较顺序，不必将杆塔座数、秒和瓦时换算为一个共同分值。第三级目标仅累计当前规划时刻至各航次结束的机载功率积分，功率以$\mathrm{W}$、时间以$\mathrm{s}$计时须除以$3600$换算为$\mathrm{Wh}$，初次规划从任务起点计起，已发生的耗电作为常数，充电电网输入另行报告，车辆牵引耗能不计入该目标。上述说明明确一次规划中的优先关系，情形集合下的推广在后文讨论。

作出这些决定，需要区分已知条件、当前状态和未来信息。地理位置、道路与障碍物描述、观测要求、设备能力及服务规则构成建模输入；车辆与无人机的位置、电池的位置与剩余能量、已完成观测和正在执行的动作构成状态输入。真实地理参考仿真能够保留任务的空间联系，但其目标几何和部分设备、服务参数采用模板或参考值，并不等于现场实测。对于尚未发生的天气变化、行驶耗时、能量消耗及观测结果，规划须依据当时可得的参考值、预测值或明确给定的不确定性描述，不能把未来实际值当作当前已知信息。方案的模型可行性也相应取决于所用输入和假设。执行中获得新信息后，后续安排须接续实际状态：已经移动到的位置、已经消耗的能量和尚未结束的动作，不能因重新分配任务而被重置。信息更新因此既可能改变方案的评价值，也可能改变后续航次、会合与充电安排的可行性。

纯组合选择与巡检联合决策共同研究资源联系下的活动组织，但承担不同的数学任务。前者在预先给定的有限候选、固定收益和资源关联中选择任务；后者还须决定路径与时序，并连接位置、能量和观测等物理状态，其完整形式包含离散变量、连续变量及非线性关系。当前容量一的组合可行域可以用两两冲突表达，超边含有较多资源本身并不构成不可约高阶耦合或求解性能优势的证据。组合实验中的选择质量和求解行为，因而不能直接迁移为巡检方案的路径、能量或会合可行性结论。要描述联合可执行性，首先需要明确哪些活动关联哪些资源，以及这些关联分别限制什么选择；超图提供了组织这类关联的表示方式，其基本对象及组合含义由此成为后续建模的起点。
